% Problem record op_299c3becfd760123. This TeX file is derived from
% database/problems_json/op_299c3becfd760123.json by scripts/sync-tex.mjs; edit the JSON record
% and rerun the script rather than editing this file.
\section{Secret key from every Bell-nonlocal behavior}
\label{sec:op_299c3becfd760123}

\paragraph{Problem.}
Does every finite-alphabet Bell-nonlocal behavior have a strictly positive
asymptotic secret-key rate against arbitrary individual nonsignalling attacks?
Let $P(a,b\mid x,y)$ be bipartite and nonsignalling, and suppose that it has no
local decomposition of the form
\begin{equation}
  P(a,b\mid x,y)
  =\int_\Lambda \mu(d\lambda)\,
    P_A(a\mid x,\lambda)P_B(b\mid y,\lambda).
  \label{eq:p30-local-decomposition}
\end{equation}
Thus Eq.~\eqref{eq:p30-local-decomposition} fails for every probability
measure $\mu$ and local response functions $P_A,P_B$.  Alice and Bob receive
independent copies of $P$; on each copy, an adversary may hold an arbitrary
nonsignalling extension.  The honest parties may choose their inputs, process
all outputs locally, and communicate publicly.  The question asks whether
some such protocol always extracts secret key at a nonzero asymptotic rate.

\subsection*{Status}

Unsolved

\subsection*{Source}

Gisin explicitly asks whether every Bell-nonlocal behavior can yield secret
key against nonsignalling individual attacks
\sourcecite{ref:p30-gisin}{Gis09}.

\subsection*{Progress}

\begin{itemize}
  \item Gisin posed the implication from Bell nonlocality to positive secret
  key specifically for nonsignalling adversaries restricted to individual
  attacks \sourcecite{ref:p30-gisin}{Gis09}.

  \item Bell nonlocality is necessary for secrecy in this adversarial model,
  and explicit nonlocal families admit positive-key protocols.  These results
  do not prove sufficiency for every behavior that violates
  Eq.~\eqref{eq:p30-local-decomposition}
  \sourcecite{ref:p30-acin}{AGM06}.

  \item Some nonlocal correlations obtained from noisy two-qubit Werner
  states have zero key rate for the broad class of standard
  device-independent protocols that announce measurements before classical
  postprocessing.  This is not an impossibility theorem for all protocols or
  for the full individual nonsignalling-adversary model
  \sourcecite{ref:p30-farkas}{FBJLA21}.

  \item Bipartite bound information exists for finite classical sources:
  there are distributions with zero key rate and positive formation cost.
  Those sources have no Bell inputs and hence do not settle the implication
  for Bell-nonlocal behaviors \sourcecite{ref:p30-pauwels}{PGR26}.
\end{itemize}

\subsection*{References}

\begin{enumerate}[label={},leftmargin=4.8em,labelsep=0.5em]
  \footnotesize
  \item[\textup{[Gis09]}]\label{ref:p30-gisin}
  N. Gisin, ``Bell Inequalities: Many Questions, a Few Answers,'' in
  W. C. Myrvold and J. Christian (eds.), \emph{Quantum Reality, Relativistic
  Causality, and Closing the Epistemic Circle}, The Western Ontario Series in
  Philosophy of Science \textbf{73}, 125--138 (Springer, 2009).
  \href{https://doi.org/10.1007/978-1-4020-9107-0_9}{doi:10.1007/978-1-4020-9107-0\_9};
  \href{https://arxiv.org/abs/quant-ph/0702021}{arXiv:quant-ph/0702021}.

  \item[\textup{[AGM06]}]\label{ref:p30-acin}
  A. Ac\'in, N. Gisin, and L. Masanes, ``From Bell's Theorem to Secure Quantum
  Key Distribution,'' \emph{Physical Review Letters} \textbf{97}, 120405
  (2006). \href{https://doi.org/10.1103/PhysRevLett.97.120405}{doi:10.1103/PhysRevLett.97.120405};
  \href{https://arxiv.org/abs/quant-ph/0510094}{arXiv:quant-ph/0510094}.

  \item[\textup{[FBJLA21]}]\label{ref:p30-farkas}
  M. Farkas, M. Balanz\'o-Juand\'o, K. \L{}ukanowski, J. Ko{\l}ody\'nski, and
  A. Ac\'in, ``Bell Nonlocality Is Not Sufficient for the Security of Standard
  Device-Independent Quantum Key Distribution Protocols,'' \emph{Physical
  Review Letters} \textbf{127}, 050503 (2021).
  \href{https://doi.org/10.1103/PhysRevLett.127.050503}{doi:10.1103/PhysRevLett.127.050503};
  \href{https://arxiv.org/abs/2103.02639}{arXiv:2103.02639}.

  \item[\textup{[PGR26]}]\label{ref:p30-pauwels}
  J. Pauwels, N. Gisin, and R. Renner, ``Bipartite Bound Information Exists,''
  arXiv:2607.25838 (2026).\newline
  \href{https://doi.org/10.48550/arXiv.2607.25838}{doi:10.48550/arXiv.2607.25838};
  \href{https://arxiv.org/abs/2607.25838}{arXiv:2607.25838}.
\end{enumerate}

\subsection*{Comment}

The unresolved implication concerns unrestricted protocols at the behavior
level under arbitrary individual nonsignalling extensions.  Impossibility
for a standard protocol class and zero-key classical sources are strictly
narrower statements.

\subsection*{Field}

Quantum Cryptography; Quantum Resource Theory

\subsection*{Topic}

Secret-key distillation; Device-independent cryptography; Bell nonlocality

\subsection*{ID}

\texttt{op\_299c3becfd760123}
